\scanpage{81}
- \mspage{IV-1015} -
fient aux éléments de $\operatorname{gr}_{\underline L}^m(N)$ appartenant à
$((R\cap\underline K^{\unsure{m}}N)+\underline L^{m+1}N)/\underline L^{m+1}N$. Il
suffira donc finalement de prouver, pour tout entier $m>0$, l'inclusion
\[
\text{(19.7.1.6)}\qquad
((R\cap\underline L^mN)+\underline L^{m+1}N)
 \cap(\underline K^mN+\underline L^{m+1}N)
 \subset(R\cap\underline K^mN)+\underline L^{m+1}N.
\]
Comme $\underline K\subset\underline L$, on vérifie aussitôt que le premier membre de (19.7.1.6)
est égal à $(R\cap(\underline K^mN+\underline L^{m+1}N))+\underline L^{m+1}N$,
si bien que tout se réduit à prouver
\[
\text{(19.7.1.7)}\qquad
\deleted{\unsure{R}}
R\cap(\underline K^mN+\underline L^{m+1}N)
 \subset(R\cap\underline K^mN)+\underline L^{m+1}N.
\]
Nous procéderons par \uline{récurrence sur m}, supposant donc (\unsure{15.7.1.7|19.7.1.7}) vérifiée
lorsqu'on remplace $m$ par un entier $m'<m$. Nous considérerons d'autre part,
pour \inserted{$m$ fixé et pour} un entier $d\geq0$, la relation
\[
(*_d)\qquad
R\cap(\underline K^mN+\underline L^{m+1}N)
 \subset(R\cap\underline K^dN\cap\underline L^mN)+\underline L^{m+1}N.
\]
\deleted{\illegible{1}} Il est clair qu'il suffit de démontrer $(*_d)$ pour $d=m$, et d'autre
part $(*_o)$ est trivialement vraie. Nous prouverons $(*_d)$ par \uline{récurrence sur d};
autrement dit, nous supposons, pour un $d\geq0$ fixé \inserted{(avec $d\leq m$)}, que $(*_d)$
est vraie, et nous voulons prouver
\[
(*_{d+1})\qquad
R\cap(\underline K^mN+\underline L^{m+1}N)
 \subset(R\cap\underline K^{d+1}N\cap\underline L^mN)+\underline L^{m+1}N.
\]
Considérons pour cela, pour un $h\geq0$, la relation
\[
(**_{d+1,h})\qquad
R\cap(\underline K^mN+\underline L^{m+1}N)
 \subset(R\cap(\underline K^{d+1}N+\underline K^d\underline L^hN)
 \cap\underline L^mN)+\underline L^{m+1}N
\]
L'hypothèse $(*_d)$ implique que $(**_{d+1,0})$ est vraie. Nous prouverons
par \uline{récurrence sur h} que $(**_{\unsure{d}+1,h})$ est vraie pour \uline{tout} $h\geq0$.
Mais \deleted{\illegible{1}}
on a le

\uline{Lemme} (19.7.1.8) .- \uline{Soient} $A$ \uline{un anneau noethérien}, $N$ \uline{un} $A$-\uline{module de}
\mspage{80}
